% ==============================================================================
% RESUMEN (español) — Página preliminar (APA 7 Nivel 1)
% Máximo 250 palabras.
% ==============================================================================

\section*{Resumen}
\addcontentsline{toc}{section}{Resumen}

El presente proyecto desarrolla un sistema de información web y móvil para optimizar la postulación y selección de estudiantes a la \enquote{Beca Comedor} de la Universidad Autónoma Gabriel René Moreno (UAGRM), gestionada por la Dirección Universitaria de Bienestar Social y Salud (DUBSS). Actualmente, la revisión de cerca de 20,000 carpetas por gestión se realiza de forma manual, lo que genera cuellos de botella, tiempos de respuesta prolongados y solicitudes que nunca llegan a ser revisadas. Además, los estudiantes carecen de un canal digital unificado para postular y consultar el estado de su solicitud.

La solución consiste en un portal web para la gestión administrativa y una aplicación móvil para los estudiantes. Su núcleo es un modelo de Inteligencia Artificial que preselecciona las solicitudes, calcula un puntaje de vulnerabilidad socioeconómica y prioriza las carpetas que el personal debe auditar. Complementariamente, un módulo de Business Intelligence ofrece tableros interactivos con indicadores sobre la preselección y el uso del comedor, que apoyan la toma de decisiones de las autoridades.

El desarrollo se conduce con el Proceso Unificado de Desarrollo de Software (PUDS), modelado con UML, y los requisitos se especifican según el estándar IEEE 830. El resultado esperado es un proceso de becas más ágil, transparente y auditable.

\indent\textit{Palabras clave: Sistema de Información, Inteligencia Artificial, Business Intelligence, PUDS, IEEE 830.}
